
\title{\LaTeX}
\section{\LaTeX}

\begin{frame}[fragile]
  \label{latex}
  \begin{block}{\centering\Large Föreläsning 3 --- \LaTeX}
    Förberedelse inför laboration 3.

    \begin{itemize}
      \ii{Ordbehandling}
      \ii{\LaTeX}
      \ii{Dokumentstruktur: dokumentklasser, omgivningar, text, stycken, listor, tabeller, \ldots}
      \ii{Programlistor}
      \ii{Matematiska formler}
      \ii{Bilder}
    \end{itemize}
  \end{block}
\end{frame}


\begin{frame}[fragile]
  \frametitle{Ordbehandling: {\sc Wysiwyg} vs. LaTeX}

  \ti{Två olika sätt att skriva text:}

  \blankline
  \begin{itemize}
    \ii{{\sc Wysiwyg} (What You See Is What You Get)}
    \begin{itemize}
      \ii{Exempel: Microsoft Word}
      \ii{Direkt återspegling av slutresultatet på skärmen}
      \ii{Inkluderar teckensnitt, storlekar, avstånd, etc.}
    \end{itemize}

    \bigskip % Lägger till extra utrymme för tydlighet
    \ii{\LaTeX\ fungerar annorlunda!}
    \begin{itemize}
      \ii{Fokus på innehåll istället för utseende}
      \ii{Skriver kommandon eller ``taggar'' för att uttrycka formatering}
      \ii{Texten \emph{kompileras} till ett formaterat dokument (oftast PDF)}
    \end{itemize}
  \end{itemize}

\end{frame}

\begin{frame}[fragile,t]
  \frametitle{\LaTeX\ - Arbetssätt och Exempel}
  \vspace{3em}

  \ti{I \LaTeX\ skriver man text och formateringskommandon i en textfil.}
  \blankline

  \begin{exlatex}
  För rätvinkliga trianglar gäller \emph{Pythagoras sats}: $a^2 + b^2 = c^2$
  \end{exlatex}

  \blankline
  \ti{Tecknet \code{\$} markerar början och slutet på matematiska uttryck. \LaTeX\ hanterar automatiskt stil och placering för matematiska element.}

\end{frame}

\begin{frame}[fragile,t]
  \frametitle{Ett större exempel}
  \vspace{3em}

  \begin{exlatex}
If $f$ is continuous on the closed interval $a \leq x \leq b$ and differentiable on the open interval $a < x < b$, then there exists a point $\xi$, $a < \xi < b$ such that:

\begin{displaymath}
  f(b) - f(a) = f'(\xi)(b - a).
\end{displaymath}
  \end{exlatex}

\end{frame}

\begin{frame}[fragile]
  \frametitle{Interaktiv Demonstration}

  \begin{center}
    \Huge Demo!

    \bigskip
    \Large Jämför {\sc Wysiwyg} med \LaTeX
  \end{center}

\end{frame}

\begin{frame}[fragile]
  \frametitle{\LaTeX-historik}
  \begin{itemize}
    \ii{Donald E. Knuth, \TeX\ från 1977}
    \ii{''Tech''—precision i typsättning}
    \ii{Leslie Lamport utvecklar \LaTeX\ på 1980-talet}
    \ii{Förenklar \TeX, fokus på innehåll}
    \ii{Bred användning inom akademi och industri}
  \end{itemize}
  \note{
    \begin{itemize}
      \item \TeX\ skapades av Donald Knuth som svar på bristerna i akademisk textproduktion.
      \item Leslie Lamport byggde vidare på detta och gjorde \TeX\ tillgängligt för en bredare publik genom \LaTeX.
      \item \LaTeX\ tillåter forskare och författare att fokusera på innehållets kvalitet snarare än dess presentation.
      \item Används för vetenskapliga dokument, böcker, och akademiska publikationer.
    \end{itemize}
  }
\end{frame}

\begin{frame}[fragile]
  \frametitle{Arbeta med \LaTeX}
  \begin{itemize}
    \ii{Skrivs i vanliga textfiler, filändelse \code{.tex}}
    \ii{Kompilera till \code{pdf} med kommandot \code{pdflatex}}
    \ii{Specialiserade \LaTeX-verktyg: \code{texmaker}, \code{TeXShop}, \code{TeXstudio}}
    \ii{VS Code med \code{LaTeX Workshop}-tillägg}
    \ii{Finns andra kompilatorer: \code{pdflatex} (vanligast), \code{XeLaTeX} (modernare, UTF-8)}
  \end{itemize}
  \note{
    \begin{itemize}
      \item \code{.tex}-filer är rena textfiler som kan redigeras med vilken textredigerare som helst.
      \item \code{pdflatex} är standardkompilatorn och fungerar bra för de flesta dokument.
      \item Integrerade verktyg som \code{texmaker}, \code{TeXShop}, och \code{TeXstudio} erbjuder enklare gränssnitt för att redigera, kompilera och visa \LaTeX-dokument.
      \item \code{VS Code} är en populär editor för många programmeringsspråk, och med \code{LaTeX Workshop}-tillägget blir den ett kraftfullt verktyg för \LaTeX.
      \item \code{XeLaTeX} erbjuder bättre stöd för moderna typografiska funktioner och fullt stöd för Unicode, vilket gör den idealisk för dokument som inkluderar flera språk eller specialtecken.
    \end{itemize}
  }
\end{frame}


\begin{frame}
  \frametitle{Fördelar med \LaTeX}
  \begin{itemize}
    \ii{Professionellt utseende}
    \ii{Automatisk numrering av rubriker, figurer, tabeller}
    \ii{\textbf{Källkoden är ren text} -- Varför är det bra? :) }
    \ii{Lätt att skapa och referera till tabeller, figurer, ekvationer}
    \ii{Automatisk innehållsförteckning}
    \ii{Bra stöd för matematiska formler}
    \ii{Gratis och plattformsoberoende, fungerar på Windows, MacOS, Linux}
    \ii{Omfattande community och resurser -- kan dock vara överväldigande och ibland stöter man på obskyra problem}
  \end{itemize}
  \note{
    \begin{itemize}
      \item \LaTeX ger dokument ett professionellt och konsekvent utseende automatiskt.
      \item Ren text-format gör det lätt att använda verktyg som Git för versionskontroll, vilket är idealiskt för samarbetsprojekt.
      \item \LaTeX-communityt erbjuder omfattande stöd och resurser, vilket gör det lätt att få hjälp och lära sig mer om avancerade funktioner.
      \item Varning för att det stora antalet paket och funktioner ibland kan kännas överväldigande och leda till att man stöter på mindre vanliga problem.
    \end{itemize}
  }
\end{frame}


\begin{frame}
  \frametitle{Dags att hoppa in i \LaTeX!}

  % "Paus"-slide innan vi går vidare till detaljer
  \begin{center}
    \Huge Nu dyker vi ner i detaljerna!
  \end{center}
\end{frame}

\begin{frame}[fragile]
  \frametitle{Använda Kommandon i LaTeX}

  \begin{itemize}
    \ii{\LaTeX-kommandon styr hur text formateras och struktureras.}
    \ii{Kan ändra typsnitt, storlek, färg, och mer.}
    \ii{De kan även organisera text i avsnitt, listor och tabeller.}
    \ii{Kommandon finns i två former: enkla kommandon och omgivningar.}
    \begin{itemize}
      \ii{Enkla kommandon: \code{\textbackslash textbf\{text\}}, \code{\textbackslash textit\{text\}}}
      \ii{Omgivningar: \code{\textbackslash begin\{itemize\}...\textbackslash end\{itemize\}}}
    \end{itemize}
  \end{itemize}
\end{frame}

\begin{frame}[fragile,t]
  \frametitle{Exempel}
  \vspace{2em}

  \begin{exlatex}
% Ex: enkelt kommando
\textbf{Fetstil} gör texten fet, \textit{kursiv stil} gör texten kursiv. Ofta används \emph{emfas} för att markera \emph{viktig text}.  
  \end{exlatex}

  \blankline
  \begin{exlatex}
% Ex: listor enumerate/itemize
\begin{enumerate}
  \item Listor hjälper till att organisera information.
  \item Varje \item skapar en ny punkt i listan.
\end{enumerate}
  \end{exlatex}
  
  \note{
    Kommandon i LaTeX skrivs med en backslash (\textbackslash) och kan ha parametrar både före och efter texten som de formaterar.
    - T.ex., \textbackslash textbf\{text\} applicerar fetstil.
    - Komplexa kommandon som \textbackslash begin\{itemize\} används för att skapa strukturerade format som listor.
  }

\end{frame}

\begin{frame}[fragile,t]
  \frametitle{Dokumentklasser}
  \vspace{3em}
  \begin{itemize}
    \ii{\emph{Dokumentklassen} styr dokumentets övergripande utseende.}
    \ii{Vanliga dokumentklasser: \code{article}, \code{report}, \code{book}, \code{letter}, \code{beamer}.}
    \ii{\code{article} används mest; \code{beamer} för presentationer.}
  \end{itemize}

  \blankline
  \begin{envi}
    \begin{lstlisting}[style=latexex]
      % Exempel på användning av dokumentklassen article
      \documentclass{article}

      % "Preamble" - inställningar och paket, t.ex.:
      \usepackage{graphicx}          % för att inkludera bilder
      \title{Enkelt LaTeX-dokument}  % sätt dokumentets titel

      \begin{document}
      \maketitle                     % skapa titelsida
      Hej världen! Detta är ett enkelt dokument.
      \end{document}
    \end{lstlisting}
  \end{envi}

  \note{
    \begin{itemize}
      \item \code{article} är idealisk för vetenskapliga artiklar, rapporter och anteckningar.
      \item \code{beamer} används för att skapa visuellt engagerande presentationer.
    \end{itemize}
  }
\end{frame}

\begin{frame}[fragile,t]
  \frametitle{Omgivningar}
  \vspace{1em}
  \begin{itemize}
    \ii{Kommandon som innehåller flera rader text kallas \emph{omgivningar}.}
    \ii{En omgivning definieras av \code{\textbackslash begin\{...\}} och \code{\textbackslash end\{...\}}.}
    \ii{Omgivningar påverkar utseendet för det innehåll de omsluter.}
  \end{itemize}

  \blankline
  \begin{exlatex}
% Exempel med omgivningar
\begin{itemize}
  \item Detta är en punktlista.
  \item Listor hjälper till att organisera information.
  \item Denna punkt har flera delar
  \begin{enumerate}
    \item Första delen
    \item Och resten
  \end{enumerate}
\end{itemize}
  \end{exlatex}

  \note{
    \begin{itemize}
      \item \code{itemize} omgivningen är bra för att skapa punktlistor.
      \item Varje \code{item} representerar en punkt i listan..
    \end{itemize}
  }
\end{frame}



\begin{frame}[fragile,t]
  \frametitle{Löpande text}
  \vspace{1em}
  \begin{itemize}
    \ii{Radslut och antal mellanslag mellan ord har ingen betydelse.}
    \ii{En eller flera blanka rader ger ett nytt stycke.}
  \end{itemize}

  \blankline
  \begin{exlatex}
Det här      är en text
 som jag  har skrivit.
Det är
en lång text med flera
     rader.

Här börjar det ett
nytt stycke i texten.
  \end{exlatex}

\end{frame}

\begin{frame}[fragile,t]
  \frametitle{Rubriker}
  \vspace{2em}
  \LaTeX\ numrerar rubriker automatiskt. Man anger en rubrik med
  \verb+\section+ eller \verb+\subsection+.

  \blankline
  \begin{exsrc}
\section{Inledning}
\section{Utförande}
\subsection{Del 1}
\subsection{Del 2}
\section{Slutsatser}
  \end{exsrc}
  \begin{exresult}
{\Large\bfseries 1 Inledning} \\[1ex]
{\Large\bfseries 2 Utförande} \\[1ex]
{\large\bfseries 2.1 Del 1} \\[1ex]
{\large\bfseries 2.2 Del 2} \\[1ex]
{\Large\bfseries 3 Slutsatser}
  \end{exresult}

\end{frame}

\begin{frame}[fragile,t]
  \frametitle{Ändra textens utseende}
  \vspace{2em}

  \begin{itemize}
    \ii{\code{\textbackslash emph\{viktigt\}} -- Betonar text}
    \ii{\code{\textbackslash texttt\{Square\}} -- Skrivmaskinstypsnitt, används ofta för kod}
  \end{itemize}

  \blankline
  
  \begin{exlatex}
Här skriver jag något \emph{viktigt}.
Och i Scala har vi skapat klassen \texttt{Square}.
I sällsynta fall kan man vilja skriva något \textbf{i fetstil}.
  \end{exlatex}
  
  \blankline
  \ti{Använd sparsamt: fetstil, lutande text, storleksändringar}
\end{frame}


\begin{frame}[fragile,t]
  \frametitle{Specialtecken}

  \begin{itemize}
    \ii{\code{\%} – Kommentar till slutet av raden}
    \ii{Tecken som kräver speciell hantering:}
  \end{itemize}

  \begin{envs}
    \begin{Code}
      \$ \% \_ \# \& \{ \} \textbackslash
    \end{Code}
  \end{envs}

  \ti{Streck, mellanrum, och punkter:}
  
  \begin{exlatex}
``pgk-kursen'' skrivs med bindestreck.
Dod-delen pågår under vecka 1--4.
Tyvärr är den inte längre\ldots

Ibland vill man göra en lång paus i en mening---för att skapa spänning.

\quad Telefon: 046--222~80~38.\\
Dagens datum: \today.
  \end{exlatex}
\end{frame}


\begin{frame}[fragile,t]
  \frametitle{Fotnoter}
  \vspace{2em}

  \ti{Fotnoter är lätta att skriva:}

  \begin{exlatexhide}
  Om man använder \LaTeX
  \footnote{uttalas
    ``lah-tekh''} så
  blir det bra. Alla rapporter
  blir automatiskt 
  % snyggt  
  professionellt
  utformade.
  \end{exlatexhide}

  \blankline
  \begin{itemize}
    \ii{Fotnoter numreras automatiskt: 1, 2, \ldots}
    \ii{Men kan även modifieras till t.ex. bokstäver eller specialtecken: *, \dag, \ddag, \S, \P, \ldots}
  \end{itemize}

\end{frame}

\begin{frame}[fragile,t]
  \frametitle{Citationstecken (1/2)}

  \ti{Använd \code{csquotes} och \code{\textbackslash enquote} i löpande text.}

\begin{exlatex}
% Svenska citat:
Hon sa \enquote{hej där} i löpande text.

% Citat i citat:
Kompisen repeterade: \enquote{Hon sa \enquote{hej där} i löpande text}.

% Raka dubbla tecken (endast kod/ASCII):
Rakt dubbelt citattecken i kod: \texttt{\textquotedbl}
\end{exlatex}

\end{frame}

\begin{frame}[fragile,t]
  \frametitle{Citationstecken (2/2)}

  \ti{Använd \code{csquotes} och \code{\textbackslash enquote} i löpande text.}

\begin{exlatex}
% Apostrof:
Typografisk apostrof: O\textquoteright{}Neill, rock\textquoteright n\textquoteright{}roll.
ASCII-apostrof (visa tecknet): \textquotesingle

% Svensk genitiv – ingen apostrof:
Annas bok; EU:s beslut (inte EU\textquoteright s).

% Engelska citat manuellt (om man inte använder csquotes):
Engelska: ``hi there'' (skrivs med \texttt{``} och \texttt{''}).
\end{exlatex}

\end{frame}


\begin{frame}[fragile,t]
  \frametitle{Listor}
  \vspace{1em}
  
  \ti{Punktlistor är enkla:}

  \begin{exlatex}
\begin{itemize}
  \item första punkten
  \item här kommer den andra
  punkten i listan
\end{itemize}
  \end{exlatex}

  \blankline
  \ti{Numrerade listor är lika enkla:}

  \begin{exlatex}
\begin{enumerate}
  \item första punkten
  \item här kommer den andra
  punkten i listan
\end{enumerate}
  \end{exlatex}

  % OBS! Pga hur detta dokument är formaterat så får listelementen inga punkter eller siffror, men i den vanliga dokumentklassen \code{\{article\}} blir det som förväntat.
  % I detta dokument används dokumentklassen \code{beamer}, och där blir numren siffror i cirklar. I den vanliga dokumentklassen \code{\{article\}} blir numren 1., 2., \ldots
\end{frame}

\begin{frame}[fragile,t]
  \frametitle{Definitioner}
  \vspace{2em}

  \ti{Man kan skapa fina definitioner:}

  \begin{exlatex}
Några klasser som vi använder:

\begin{description}
  \item[Scanner] Inläsning från tangentbordet
  \item[Random] Slumptal
  \item[File] För att läsa och skriva filer
\end{description}
  \end{exlatex}

\end{frame}

\begin{frame}[fragile,t]
  \frametitle{Tabeller}
  \vspace{2em}
  
  \ti{Tabeller är enkla att skapa:}
  \begin{itemize}
    \ii{Använd \code{tabular}-omgivningen}
    \ii{Specificera antal kolumner och hur texten ska justeras}
    \ii{Använd \code{\&} för att skilja kolumner och \code{\textbackslash\textbackslash} för att skilja rader}
  \end{itemize}

  \begin{exlatex}
\begin{tabular}{lcr}
  Produkt & Typ    & Pris    \\
  \hline
  Skruvar & stora  & 0.18~kr \\
  Muttrar & M16    & 0.38~kr \\
  Spikar  & 12~tum & 0.12~kr
\end{tabular}
  \end{exlatex}

\end{frame}

\begin{frame}[fragile,t]
  \frametitle{Flytande tabeller}
  \vspace{2em}

  \ti{Med en \texttt{table}-omgivning skapar man en \emph{flytande tabell}, som \LaTeX\ själv placerar där det är lämpligt.}

  \ti{Man kan ange sin preferens för placering med \texttt{h}~(here), \texttt{t}~(top), \texttt{b}~(bottom), \texttt{p}~(page).}

  \begin{exlatex}
\begin{table}[t]
  \begin{tabular}{|l|c|r|}
    Produkt & Typ    & Pris    \\
    \hline
    Skruvar & stora  & 0.18~kr \\
    Muttrar & M16    & 0.38~kr \\
    Spikar  & 12~tum & 0.12~kr
  \end{tabular}
  \caption{Våra produkter}
\end{table}
  \end{exlatex}
\end{frame}

\begin{frame}[fragile,t]
  \frametitle{Att referera till etiketter}
  \vspace{1em}

  \ti{Med \emph{etiketter} (label) kan man referera till tabeller från texten:}

  \ti{(Figurer hanteras likadant som tabeller, i en \texttt{figure}-omgivning.)}

  \begin{exlatex}
\begin{table}
  \begin{tabular}{lcr}
    Produkt & Typ    & Pris    \\
    \hline
    Skruvar & stora  & 0.18~kr \\
    Muttrar & M16    & 0.38~kr \\
    Spikar  & 12~tum & 0.12~kr
  \end{tabular}
  \caption{Våra produkter}
  \label{produkter}
\end{table}
Senare i texten kan vi referera till vår tabell \ref{produkter}, som visar våra produkter.

Senare i texten kan vi referera till vår tabell~\ref{produkter}, som visar våra produkter.
  \end{exlatex}

\end{frame}

\begin{frame}[fragile,t]
  \frametitle{Programlistor}
  \vspace{.5em}

  \begin{itemize}
    \ii{För programkod kan man använda \texttt{verbatim}-omgivningen.}
    \ii{Den visar texten precis som den är skriven.}
    \ii{Med paketet \texttt{fancyvrb} kan man också använda \texttt{\textbackslash VerbatimInput}-kommandot för att inkludera kod från en fil.}
    \ii{\textbf{Notera} att \texttt{verbatim} inte hanterar tabulatortecken korrekt!}
  \end{itemize}

  \halfblankline
  \begin{exsrc}
\usepackage{fancyvrb}
\fvset{tabsize=4, fontsize=\small}

\VerbatimInput{Point.scala}
\end{exsrc}
\begin{exresult}%
{\small
\begin{verbatim}
class Point(val x: Int, 
            val y: Int) {
  
  def move(dx: Int, 
            dy: Int): Point = 
    new Point(x + dx, y + dy)
}
\end{verbatim}
}  
  \end{exresult}
\end{frame}

\begin{frame}[fragile,t]
  \frametitle{Programlistor med \texttt{listings}}
  \vspace{.5em}

  \begin{itemize}
    \ii{Paketet \texttt{listings} är ett bättre alternativ.}
    \ii{Stöd för syntaxfärgning och flera programmeringsspråk.}
    \ii{Många anpassningsmöjligheter: teckensnitt, färger, tabstorlek, \ldots}
  \end{itemize}

  \begin{exsrc}
\lstset{
  language=Scala,
  basicstyle=\ttfamily\small,
  % fler inställningar ...
}

\begin{lstlisting}
  class Point(val x: Int, 
    val y: Int) {

      def move(dx: Int, 
        dy: Int): Point = 
      new Point(x + dx, y + dy)
  }
\end{lstlisting}
  \end{exsrc}
  \begin{exresult}
\begin{lstlisting}[style=scalaex]
class Point(val x: Int, 
  val y: Int) {

    def move(dx: Int, 
      dy: Int): Point = 
    new Point(x + dx, y + dy)
}
  \end{lstlisting}
  \end{exresult}
\end{frame}


\begin{frame}[fragile]
  \frametitle{Öka eller minska avstånd}

  \begin{itemize}
    \ii{Vertikalt avstånd: \texttt{\textbackslash vspace\{längd\}}}
    \ii{Längder kan anges i cm, mm, in, pt, \ldots}
    \ii{Negativ längd minskar avståndet.}
    \ii{Snabbkommandon: \code{\textbackslash smallskip}, \code{\textbackslash medskip}, \code{\textbackslash bigskip}}
    \ii{Horisontellt avstånd: \texttt{\textbackslash hspace\{längd\}}}
  \end{itemize}
\end{frame}


\begin{frame}[fragile]
  \frametitle{Matematiska formler}
  \ti{\LaTeX\ är mycket bra på att formatera matematisk text.}

  \blankline
  \begin{itemize}
    \ii{I löpande text: \code{\$ \dots\ \$} eller \code{\textbackslash( \dots\ \textbackslash)} (mer robust i vissa sammanhang.)}
    \ii{På egen rad: \texttt{\textbackslash begin\{displaymath\} \dots\ \textbackslash end\{displaymath\}}\\
      eller \code{\textbackslash [ \dots\ \textbackslash ]} }
    \ii{Numrerad formel: \code{\textbackslash begin\{equation\} \dots\ \textbackslash end\{equation\}}}
    \ii{Etikett och referens: \code{\textbackslash label} och \code{\textbackslash ref}}
  \end{itemize}

\end{frame}


\begin{frame}[fragile,t]
  \frametitle{Enkla formler}
  \vspace{2em}

  \begin{exlatex}
Formeln \(x=3y-2\) står inne i texten. Däremot står
\begin{displaymath}
  x_1 = 3y - 2
\end{displaymath}
för sig själv, precis som
\[x_2 = -0.6e^{2-y} + \frac{9.3}{\pi}\]
och detta är en numrerad ekvation:
\begin{equation}
  x_3 = 2y^2 + 7y - 2a + 3
  \label{xochy}
\end{equation}

I ekvation~\ref{xochy} fann
vi att \ldots
  \end{exlatex}
\end{frame}

\begin{frame}[fragile,t]
  \frametitle{Symboler, index}
  \vspace{2em}

  \begin{exlatex}
\begin{displaymath}
  \alpha \leq \pi \approx 3.141592654
\end{displaymath}
  \end{exlatex}

  \vspace{1cm}

  \begin{exlatex}
\[
  x_{k+1} =
  x_{k} -
  \frac{f(x_{k})}{f'(x_{k})}
\]
  \end{exlatex}
\end{frame}

\begin{frame}[fragile,t]
  \frametitle{Exponenter, rötter}
  \vspace{2em}

  \begin{exlatex}
\[
  e^x = 1+x+x^2/2!+x^3/3!+\cdots
\]
  \end{exlatex}

  \vspace{1cm}

  \begin{exlatex}
\[
  x_{1,2}=\frac{p}{2}\pm
  \sqrt{\frac{p^2}{4}-q}
\]
  \end{exlatex}

\end{frame}



\begin{frame}[fragile,t]
  \frametitle{Integraler, summor}
  \vspace{2em}

  \begin{exlatex}
\[
  \int_{-\infty}^{\infty}
  e^{-x^2} dx
\]
  \end{exlatex}

  \vspace{1cm}

  \begin{exlatex}
\[
  \sum_{k=1}^n\frac{1}{a_k}
\]
  \end{exlatex}

\end{frame}

\begin{frame}[fragile,t]
  \frametitle{Funktioner}
  \vspace{2em}

  \begin{exlatex}
\[
  \sin^2 x + \cos^2 x = 1
\]
  \end{exlatex}
\end{frame}

\begin{frame}[fragile,t]
  \frametitle{Matriser, parenteser}
  \vspace{2em}

  \begin{exlatex}
\[
A=\left(
\begin{array}{cccc}
a_{11} & a_{12} & \cdots & a_{1n} \\
a_{21} & a_{22} & \cdots & a_{2n} \\
\vdots & \vdots & \ddots & \vdots \\
a_{n1} & a_{n2} & \cdots & a_{nn} \\
\end{array}
\right)
\]
  \end{exlatex}
\end{frame}


\begin{frame}[fragile,t]
  \frametitle{Bilder}
  \vspace{2em}

  \ti{Inkludera bilder i format: \code{pdf}, \code{jpeg}, \code{png}}

  \ti{Använd paketet \texttt{graphicx} (eller \code{graphics})}

  \blankline

  \vspace{2mm}
  \begin{exlatex}
  \includegraphics[height=30mm]
    {images/bild.pdf}
  \end{exlatex}
%   \begin{exsrc}
% \includegraphics[height=40mm]
%     {images/bild.pdf}
%   \end{exsrc}
%   \begin{exresult}
% \includegraphics[height=30mm]{images/enkelmodell.pdf}
%   \end{exresult}

\end{frame}

\begin{frame}[fragile,t]
  \frametitle{Bilder med figurmiljö}
  \vspace{1em}

  \ti{Slå in bilder i en \texttt{figure}-omgivningen, för att lägga till en bildtext och en etikett.}

    \begin{exlatex}
\begin{figure}
  \centering
  \includegraphics[height=30mm]
      {images/enkelmodell.pdf}
  \caption{Exempelbild}
  \label{fig:exempelbild}
\end{figure}
Från text kan vi nu referera till figur~\ref{fig:exempelbild}.
    \end{exlatex}

\end{frame}


\begin{frame}[fragile]
  \frametitle{Egna kommandon}
  Man kan lätt definiera egna kommandon, till exempel ett kortare namn
  för en text som man använder ofta. Kommandon kan ha parametrar.

  \vspace{10mm}

  \begin{exlatex}
\newcommand{\scala}[1]
  {\texttt{\textcolor{red}{#1}}}

Klasser: \scala{Random},
\scala{Scanner}, \scala{FileUtils} och
\scala{PrintStream}.
  \end{exlatex}

  \vspace{10mm}

  Man kan definiera om existerande kommandon med
  \verb+\renewcommand+. Det kan ställa till förvirring, så
  gör inte det.
\end{frame}

\begin{frame}[fragile]
  \frametitle{\LaTeX\ på egen dator}
  
  \ti{En sammanfattning av LaTeX-installationer finns på \url{www.latex-project.org}, sidan Getting LaTeX.}

  \blankline
  \begin{description}
    \di{Linux}{Installera LaTeX via den vanliga pakethanteraren (t.ex. \texttt{apt}).}
    \di{Mac}{Använd \texttt{MacTeX} (baserad på \texttt{TeXLive}, uppdateras årligen).}
    \di{Windows}{Rekommenderat: \texttt{TeXLive} (\url{www.tug.org/texlive/}) för en komplett LaTeX-distribution.}
  \end{description}
  
  \halfblankline
  \begin{itemize}
    \ii{För redigering, använd moderna verktyg som t.ex. \texttt{TeXstudio}, \texttt{TeXShop} (för Mac), eller \texttt{VS Code} med \texttt{LaTeX Workshop}-tillägget.}
    \ii{\textbf{Notera}, ni får lov att använda Overleaf om ni föredrar det!}
    \ii{Lunds Universitet erbjuder nu t.o.m. gratis premiumkonto på Overleaf för studenter och anställda vid LU. Se mer här: \url{https://luservicedesk.service-now.com/support?id=kb_article_swe&sys_id=1be990302b6ba6900002fa97d891bf77}}
  \end{itemize}

\end{frame}
